% GENERATED FILE. Edit canonical lessons, then run npm run generate:books.
% canonical-source-hash: fnv1a64:ca4319edaf1fe4a2
% canonical-lessons: BN-C270-dur, BN-W270-letter-uu-kar, BN-C270-kibhabe-jabo, BN-R270-asking-the-way

\chapter{Which Way, and How Far}
\label{ch:bn-which-way-and-how-far}
\hlchaptermodality{270}

\begin{chapteropening}
\textbf{By the end of this chapter:} \emph{I can ask the way (kībhābe jābô? kôtô dūr?), give it (sojā, bā̃dike, ḍāndike), and read and write the sign \bn{ূ} in \bn{দূর}.}

Complete BN-R270-asking-the-way: ask the way and give it.
\end{chapteropening}

\section[\texorpdfstring{\bn{দূর} (dūr)}{দূর (dūr)}]{\bn{দূর} (dūr) — far; a distance}

\label{lesson:BN-C270-dur}



\begin{quote}
Before the new one: say \textquotedblleft{}so we stayed at home\textquotedblright{}. (\emph{tāi āmrā bāṛite thāklām}.) Then \textbf{\bn{কাছে}}, \textquotedblleft{}near\textquotedblright{}.
\end{quote}

\subsection*{You'll want to know: \bn{দূর}}
\textbf{\bn{দূর}} — \emph{dūr} — \textquotedblleft{}far; a distance\textquotedblright{}.

\begin{quote}
\textbf{\bn{বাজার} \bn{আমাদের} \bn{বাড়ি} \bn{থেকে} \bn{দূরে।}}
\end{quote}

\begin{quote}
\emph{bājār āmāder bāṛi theke dūre.} — \textquotedblleft{}The market is far from our house.\textquotedblright{}
\end{quote}

\textbf{\bn{দূরে}}, with \textbf{-\bn{এ}}, is \textquotedblleft{}far away, at a distance\textquotedblright{}, the opposite of \textbf{\bn{কাছে}}. \textbf{\bn{থেকে}} gives the point it is measured from. \textbf{\bn{দূর} \bn{থেকে}} is \textquotedblleft{}from far away\textquotedblright{}: \textbf{\bn{নদীটা} \bn{দূর} \bn{থেকে} \bn{দেখলাম}} (\emph{nôdīṭā dūr theke dekhlām}), \textquotedblleft{}I saw the river from far away\textquotedblright{}.

\begin{cousinweb}
\textbf{\bn{দূর}} (\emph{dūr}) is a tatsama: Sanskrit \emph{dūra}, \textquotedblleft{}distant\textquotedblright{}, and as a noun \textquotedblleft{}distance\textquotedblright{} (CDIAL 6495). CDIAL lists the Bengali word under it as \emph{dur}: Bengali does not lengthen the vowel, and the long \textbf{\bn{ূ}} is the Sanskrit spelling kept.
\end{cousinweb}

\subsection*{Your turn}
Say these aloud:

\begin{itemize}
\raggedright
  \item \emph{dūr}
  \item \textquotedblleft{}The market is far from our house\textquotedblright{} — \emph{bājār āmāder bāṛi theke dūre}
  \item \textquotedblleft{}The station is near\textquotedblright{} — \emph{sṭeshôn kāchhe}
\end{itemize}

\begin{tcolorbox}[breakable,colback=teal!4,colframe=teal!35!black,title={Before you move on}]
What does \textbf{\bn{দূর}} mean? (\textquotedblleft{}Far; a distance\textquotedblright{}.) \textquotedblleft{}Far away\textquotedblright{}? (\textbf{\bn{দূরে}}.) Its opposite? (\textbf{\bn{কাছে}}.)
\end{tcolorbox}
\section[\texorpdfstring{\bn{ূ} (-ū)}{ূ (-ū)}]{\bn{ূ} (-ū) — the sign under the \bn{দ} of \bn{দূর}}

\label{lesson:BN-W270-letter-uu-kar}



\begin{quote}
Before the new one: say \textbf{\bn{দূর}}, \textquotedblleft{}far\textquotedblright{}. Then write \textbf{\bn{তু}} and \textbf{\bn{উ}}: the \emph{u} sign and the \emph{u} letter you know.
\end{quote}

\subsection*{Script you'll notice: \bn{ূ}}
\textbf{\bn{ূ}} — \emph{-ū}, the long \emph{u}, written under its consonant: \textbf{\bn{দূর}}. Bengali says it the same as \textbf{\bn{ু}}; the long sign is a spelling, mostly in words from Sanskrit, as \textbf{\bn{ী}} is beside \textbf{\bn{ি}}.

Set the print of \textbf{\bn{দূ}} beside \textbf{\bn{দু}}. Under \textbf{\bn{দু}} hangs a small closed loop with a short slanting tail at its lower right. Under \textbf{\bn{দূ}} hangs a rounded hook that is open on the right, like a C, with a short slanting tail at its lower end.

\subsection*{Writing: \bn{ূ} — copy the printed shape}
Put your pen beside \textbf{\bn{দূ}} and copy the shape you can see, slowly and large: the \textbf{\bn{দ}}, then the open hook under it, then its tail. Check the copy against the print: open, not closed like \textbf{\bn{ু}}.

\begin{quote}
This book does not yet tell you where to start \bn{ূ} or which way the pen goes. That is not written down here until it can be written down with a source. Copying the printed shape needs no such source, and it is how the shape becomes yours in the meantime.
\end{quote}

\subsection*{Your turn}
\begin{itemize}
\raggedright
  \item \emph{Find:} \bn{ূ} under the \bn{দ} of \textbf{\bn{দূর}}
  \item \emph{Write it:} \textbf{\bn{দু}}, then \textbf{\bn{দূ}}, from the print
  \item \emph{Write it:} \textbf{\bn{দূর}}, then \textbf{\bn{দূরে}}
  \item \emph{Write it:} \textbf{\bn{উচিত}}, from memory
\end{itemize}

\begin{tcolorbox}[breakable,colback=teal!4,colframe=teal!35!black,title={Before you move on}]
Which sign is under the \textbf{\bn{দ}} of \textbf{\bn{দূর}}? (\textbf{\bn{ূ}}.) How is it unlike \textbf{\bn{ু}}? (Its hook is open on the right.) Is it said longer? (No: the same \emph{u}.)
\end{tcolorbox}
\section[\texorpdfstring{… \bn{কীভাবে} \bn{যাব}? \bn{কত} … (… kībhābe jābô? …)}{… কীভাবে যাব? কত … (… kībhābe jābô? …)}]{… \bn{কীভাবে} \bn{যাব}? \bn{কত} \bn{দূর}? (kībhābe jābô? kôtô dūr?) — asking the way}

\label{lesson:BN-C270-kibhabe-jabo}



\begin{quote}
Say \textquotedblleft{}The market is far from our house\textquotedblright{}. (\emph{bājār āmāder bāṛi theke dūre}.) Then \textquotedblleft{}I will go\textquotedblright{}. (\emph{jābô}.)
\end{quote}

\subsection*{You'll want to know: asking}
\begin{quote}
\textbf{\bn{স্টেশনে} \bn{কীভাবে} \bn{যাব}?}
\end{quote}

\begin{quote}
\emph{sṭeshône kībhābe jābô?} — \textquotedblleft{}How do I get to the station?\textquotedblright{}
\end{quote}

\textbf{\bn{কীভাবে}} (\emph{kībhābe}), \textquotedblleft{}how\textquotedblright{}, is \textbf{\bn{কী}}, \textquotedblleft{}what\textquotedblright{}, with \textbf{\bn{ভাবে}}, \textquotedblleft{}in the manner\textquotedblright{}: \textquotedblleft{}how shall I go?\textquotedblright{}. \textbf{\bn{কত}} (\emph{kôtô}), \textquotedblleft{}how much\textquotedblright{}, with \textbf{\bn{দূর}} asks the distance:

\begin{quote}
\textbf{\bn{স্টেশন} \bn{কত} \bn{দূর}?}
\end{quote}

\begin{quote}
\emph{sṭeshôn kôtô dūr?} — \textquotedblleft{}How far is the station?\textquotedblright{}
\end{quote}

To ask a stranger, start with \textbf{\bn{মাফ} \bn{করবেন}}, \textquotedblleft{}excuse me\textquotedblright{}.

\subsection*{You'll want to know: giving the way}
The answer often uses the polite future, \textquotedblleft{}you will go\textquotedblright{}, with \textbf{\bn{সোজা}}, \textquotedblleft{}straight on\textquotedblright{}, \textbf{\bn{বাঁদিকে}} (\emph{bā̃dike}), \textquotedblleft{}to the left\textquotedblright{}, and \textbf{\bn{ডানদিকে}} (\emph{ḍāndike}), \textquotedblleft{}to the right\textquotedblright{}: \textbf{\bn{বাঁ}}, \textquotedblleft{}left\textquotedblright{}, \textbf{\bn{ডান}}, \textquotedblleft{}right\textquotedblright{}, with \textbf{\bn{দিক}}, \textquotedblleft{}a side\textquotedblright{}, and \textbf{-\bn{এ}}.

\begin{quote}
\textbf{\bn{সোজা} \bn{গিয়ে} \bn{মোড়} \bn{থেকে} \bn{ডানদিকে} \bn{যাবেন।} \bn{স্টেশন} \bn{বাজারের} \bn{পাশে।}}
\end{quote}

\begin{quote}
\emph{sojā giye moṛ theke ḍāndike jāben. sṭeshôn bājārer pāshe.} — \textquotedblleft{}Go straight, then turn right at the corner. The station is next to the market.\textquotedblright{}
\end{quote}

And the distance: \textbf{\bn{দশ} \bn{মিনিট} \bn{হাঁটতে} \bn{হবে}} (\emph{dôsh miniṭ hā̃ṭte hôbe}), \textquotedblleft{}it's a ten-minute walk\textquotedblright{}.

\subsection*{Your turn}
Say these aloud:

\begin{itemize}
\raggedright
  \item \textquotedblleft{}Excuse me, how do I get to the market?\textquotedblright{} — \emph{māph kôrben, bājāre kībhābe jābô?}
  \item \textquotedblleft{}Go straight, then turn left\textquotedblright{} — \emph{sojā giye bā̃dike jāben}
  \item \textquotedblleft{}It's a ten-minute walk\textquotedblright{} — \emph{dôsh miniṭ hā̃ṭte hôbe}
\end{itemize}

\begin{tcolorbox}[breakable,colback=teal!4,colframe=teal!35!black,title={Before you move on}]
\textquotedblleft{}How far?\textquotedblright{} (\textbf{\bn{কত} \bn{দূর}?}) \textquotedblleft{}To the left\textquotedblright{}? (\textbf{\bn{বাঁদিকে}}.) \textquotedblleft{}To the right\textquotedblright{}? (\textbf{\bn{ডানদিকে}}.)
\end{tcolorbox}
\section[(review)]{Review — which way, and how far}

\label{lesson:BN-R270-asking-the-way}



\begin{quote}
Ask \textquotedblleft{}How far is the station?\textquotedblright{} (\emph{sṭeshôn kôtô dūr?})
\end{quote}

\subsection*{Your turn}
\begin{itemize}
\raggedright
  \item \emph{Say it:} Excuse me, how do I get to the hospital?
  \item \emph{Say it:} Go straight, then turn right at the corner.
  \item \emph{Say it:} The hotel is next to the station. It isn't far.
  \item \emph{Say it:} How far is the river? It's a ten-minute walk.
  \item \emph{Say it:} I don't think the market is far, but you ought to take a bus.
  \item \emph{Write it:} \textbf{\bn{দূর}}, then \textbf{\bn{উচিত}}
\end{itemize}

\begin{tcolorbox}[breakable,colback=teal!4,colframe=teal!35!black,title={Before you move on}]
\textquotedblleft{}How do I get to the station?\textquotedblright{} (\textbf{\bn{স্টেশনে} \bn{কীভাবে} \bn{যাব}?}) \textquotedblleft{}Far away\textquotedblright{}? (\textbf{\bn{দূরে}}.) Which sign is under the \textbf{\bn{দ}} of \textbf{\bn{দূর}}? (\textbf{\bn{ূ}}.)
\end{tcolorbox}
